\begin{table}[!ht]
\centering
\caption{Decomposição do gap salarial racial por mediação contextual e ocupacional
(modelos HLM de três níveis, PNAD Contínua 2016--2025, população completa). Cada linha
acrescenta um bloco de controles ao anterior; à medida que se adiciona contexto e ocupação,
o gap encolhe e a mediação acumulada cresce.}
\label{tab:mediacao}
\begin{tabular}{lrrr}
\toprule
Modelo (controles acumulados) & $\beta_{\text{negro}}$ & Gap (\%) & Mediação acum. (\%) \\
\midrule
Individual (educ., idade, sexo) & −0,2123 & −19,1 & --- \\
+ contexto da UPA (vizinhança) & −0,1009 & −9,6 & 52,5 \\
+ efeito de UF (estado) & −0,1009 & −9,6 & 52,5 \\
+ ocupação e formalidade & −0,0641 & −6,2 & 69,8 \\
\bottomrule
\end{tabular}
\par\smallskip
\footnotesize\emph{Como ler:} cada linha adiciona controles à anterior. $\beta_{\text{negro}}$
é a penalidade racial em log-rendimento (mais próximo de zero = menor gap); ``Gap (\%)'' é a
penalidade em \% de renda; ``Mediação acum.''\ é a fração do gap bruto (M1) já explicada pelos
controles. Do M1 ao M4 o gap cai de $-19{,}1\%$ para $-6{,}2\%$: $69{,}8\%$ do gap é mediado por
onde a pessoa mora e qual ocupação acessa --- e o residual de $-6{,}2\%$ persiste após todos os
controles observáveis.
\end{table}
